\documentclass[12pt]{article}
\usepackage[utf8]{inputenc}
\usepackage[T1]{fontenc}
\usepackage{hyperref}
\title{Awyu, Asue Language Resources for K-12 Curriculum Development}
\author{Compiled by Claude Sonnet 5.5 (claude-sonnet-5-5) for Andr\'as Kornai}
\date{}
\begin{document}
\maketitle

\section{Introduction}
Asue Awyu (ISO 639-3 \texttt{psa}) is an endangered Papuan language of Indonesian Papua in the Greater Awyu family, with about 6,500 speakers according to the task description. It is not known to be taught in schools. The searches found nothing specific to Asue.

\section{Primary Educational Resources}
None found.

\section{Linguistic and Reference Materials}
\begin{itemize}
\item Lourens de Vries, \emph{The Greater Awyu Languages of West Papua}, a description of the family. Whether it covers Asue in detail was not confirmed from the search results: \url{https://vufind.lib.tsu.ru/Record/libtsu.1013453/Details}
\end{itemize}

\section{Religious and Community Resources}
None found.

\section{Digital and Audio Resources}
None found.

\section{Curriculum Development Context}
There is essentially no Asue-specific material in what was found. Any curriculum would start from documentation and orthography work with the community, with Indonesian as the school language.
\end{document}